\guidepart{amber}{Stay safe}{Stay safe}
  {Passwords, keys and host checks, minus the jargon headache.}
  {\partitem{\iLock}{Passwords and the vault}
   \partitem{\iShield}{Host keys: is it really that server?}
   \partitem{\iKey}{SSH keys}}

% =========================================================================================
\chapter{Passwords and the vault}\label{ch:secrets}
\chapmeta{5 minutes}{You'll know exactly where your passwords live and who can read them}

\begin{tldr}
Passwords go in your OS keychain. If there isn't one, they go in an encrypted vault protected by a master password.
For each password you choose: save it, remember it until you quit, or ask every time.
\end{tldr}

\section{Where secrets live}

\begin{center}
\begin{tikzpicture}[
  box/.style={rounded corners=6pt, fill=#1!12, draw=#1!60, line width=0.6pt, minimum height=13mm, text width=40mm,
    align=center, font=\small},
  arr/.style={-latex, line width=0.8pt, color=muted}]
  \node[box=accent] (need) {\heavy A session needs\\a password};
  \node[box=green, right=12mm of need] (kc) {\heavy OS keychain\\{\footnotesize Keychain · Credential Manager · Secret Service}};
  \node[box=amber, right=12mm of kc] (vault) {\heavy Encrypted vault\\{\footnotesize AES-256-GCM, Argon2id}};
  \node[box=rose, below=9mm of kc] (ask) {\heavy Ask you\\{\footnotesize then carry on}};
  \draw[arr] (need) -- (kc);
  \draw[arr] (kc) -- node[above, font=\scriptsize]{not there?} (vault);
  \draw[arr] (vault) |- node[pos=0.25, right, font=\scriptsize]{still not there?} (ask);
\end{tikzpicture}
\end{center}

\begin{description}
  \item[OS keychain first] macOS Keychain, Windows Credential Manager, or the Linux Secret Service (GNOME Keyring,
        KWallet and friends).
  \item[Encrypted vault as backup] If there's no keychain, or it refuses, secrets go into Nexus Terminal's own
        vault: AES-256-GCM with a key derived from your master password by Argon2id using 128\,MiB of memory. In
        plain English: very, very expensive to brute-force.
\end{description}

\section{Setting up the vault}

\menu{Settings\then Security\then Initialise Secret Store}. Choose a master password; the strength meter goes from
\proto[rose]{Weak} through \proto[amber]{Fair} and \proto[accent]{Good} to \proto[green]{Strong}. Unlock and lock it
from the same page, from the palette, or by clicking the vault badge in the status bar. It locks itself when you quit.

\crop{settings-security.png}{581}{253}{576}{285}{0.8\linewidth}{Security settings on a machine without an OS keychain: the vault does the job instead.}

\begin{danger}[There's no ``forgot password'' link]
Lose the master password and the secrets in the vault are gone for good. That's the point of encryption. Put it in
your password manager.
\end{danger}

\section{Three ways to handle each password}

Every password, passphrase, token or TOTP secret has its own setting:

\begin{center}
\begin{tabular}{@{}p{0.3\linewidth}p{0.3\linewidth}p{0.3\linewidth}@{}}
  \begin{tcolorbox}[colback=green!8, boxrule=0pt, arc=6pt, fontupper=\small, height=33mm]
    {\heavy\color{green!60!black}\iLock\ Save securely}\par\smallskip Stored in the keychain or vault and used automatically. Most convenient.
  \end{tcolorbox} &
  \begin{tcolorbox}[colback=accent!8, boxrule=0pt, arc=6pt, fontupper=\small, height=33mm]
    {\heavy\color{accent!80!black}\iClock\ Remember for this session}\par\smallskip Asked once, kept in memory, never written to disk.
  \end{tcolorbox} &
  \begin{tcolorbox}[colback=amber!10, boxrule=0pt, arc=6pt, fontupper=\small, height=33mm]
    {\heavy\color{amber!75!black}\iDetective\ Ask every time}\par\smallskip Asked on every single connection. Most paranoid.
  \end{tcolorbox}
\end{tabular}
\end{center}

When Nexus Terminal needs a secret it doesn't have, it asks. Press \kbd{Esc} or \ui{Cancel} to give up on that
connection.

\section{Seeing and changing what's saved}

\ui{Stored credentials} on the Security page lists every saved password, passphrase and token by profile, and whether
it's in the keychain or the vault. \ui{Forget} deletes one. \ui{Change Master Password…} re-encrypts the whole vault
under a new password in one go.

\begin{sidequest}[What gets written down?]
Connections, disconnections, profile changes, credential access, imports and exports go into an audit log in the
database. Secrets themselves never do. TOTP secrets live in the keychain or vault too, never in the profile.
\end{sidequest}

% =========================================================================================
\chapter{Host keys: is it really that server?}\label{ch:hostkeys}
\chapmeta{3 minutes}{You'll understand the scary fingerprint dialog and what to do about it}

\begin{tldr}
Every SSH server (and every RDP server, via its certificate) has a fingerprint. Nexus Terminal remembers it the first time and shouts if it ever changes.
A changed fingerprint means a reinstalled server, or someone in the middle. Find out which before carrying on.
\end{tldr}

Pick a policy per SSH profile under \ui{Host Verification}:

\renewcommand\arraystretch{1.5}
\begin{tabularx}{\linewidth}{@{}>{\leavevmode\heavy\color{ink}}p{48mm}X@{}}
  \toprule
  Strict (known hosts only)   & Only connects to servers you've already trusted.\\
  Ask on new or changed keys  & Shows the fingerprint and asks. New: \ui{Accept \& Save} or \ui{Reject}. Changed:
                                \ui{Reject}, \ui{Accept Once} or \ui{Accept \& Update}. The default.\\
  Accept new, reject changed  & Trusts first-time servers silently, still refuses changed keys.\\
  Insecure (accept anything)  & Exactly what it says. Lab use only.\\
  \bottomrule
\end{tabularx}

Under \emph{Strict} and \emph{Accept new}, a key that differs from the saved one is rejected outright.

\crop{settings-knownhosts.png}{581}{253}{576}{285}{0.8\linewidth}{Known hosts. Click a fingerprint to copy it; \ui{Remove} to forget a server.}

\begin{quickwin}
Nexus Terminal keeps its own list, but \ui{Import \textasciitilde/.ssh/known\_hosts} brings in everything your
command-line \code{ssh} already trusts (hashed names can't be read back, so those are skipped). \ui{Export…} goes the
other way, in OpenSSH format.
\end{quickwin}

\section{Certificate authorities}

Big estates sign host keys with an SSH certificate authority instead of handing out fingerprints. \ui{Trust a CA…}
adds the CA's public key for a host pattern in OpenSSH style (\code{*.corp.example}, \code{[jump.example]:2222},
\code{!old.corp.example}); \code{@cert-authority} lines in an imported \code{known\_hosts} do the same. A server the CA
covers must then show a certificate the CA signed, naming that host and in date. It's trusted without a prompt, and
anything else is refused.

% =========================================================================================
\chapter{SSH keys}\label{ch:keys}
\chapmeta{2 minutes}{You'll have a key pair without memorising ssh-keygen flags}

\menu{Session\then SSH key manager} (or the palette) lists the key pairs in \code{\textasciitilde/.ssh} with their type. Copy a
public key to the clipboard (ready to paste into GitHub or \code{authorized\_keys}), delete a pair, or generate a new
\proto[teal]{Ed25519}, \proto[accent]{RSA} (2048 or 4096 bits) or \proto[violet]{ECDSA} key with an optional passphrase.

\crop{keys.png}{725}{523}{725}{546}{0.75\linewidth}{The SSH key manager.}

\begin{protip}
Not sure which type? Pick Ed25519. It's small, fast and secure. Keys are generated inside Nexus Terminal and
saved in the usual OpenSSH format, so \code{ssh} and \code{git} can use them too.
\end{protip}
